\documentclass[12pt]{article}
\usepackage[utf8]{inputenc}
\usepackage[T1]{fontenc}
\usepackage{hyperref}
\title{San Juan Atzingo Popoloca Language Resources for K-12 Curriculum Development}
\author{Compiled by Claude Sonnet 5.5 (claude-sonnet-5-5) for Andr\'as Kornai}
\date{}
\begin{document}
\maketitle

\section{Introduction}
San Juan Atzingo Popoloca (ISO 639-3 poe; Glottocode sanj1285) is a Popolocan language of the Otomanguean family spoken in Puebla, Mexico. Population figures differ between sources (about 6,900 to 14,500). It is also called Eastern Popoloca, Southern Popoloca or Ngiba. Searches found Scripture, audio and older literacy materials but no current school curriculum.

\section{Primary Educational Resources}
Older SIL literacy primers (1973--1989) are digitised at the repository of El Colegio de M\'exico, and the notes there say that their orthography is no longer current. The individual item pages were not confirmed from the search results, so none is listed separately here.

\section{Linguistic and Reference Materials}
\begin{itemize}
\item Glottolog entry: \url{https://glottolog.org/resource/languoid/id/sanj1285}
\item WALS language page, which lists 1995 reference works (Kalstrom Dolson et al.; Krumholz et al.) as sources: \url{https://wals.info/languoid/lect/wals_code_psj}
\item Wikitongues language page: \url{https://wikitongues.org/languages/poe/}
\end{itemize}

\section{Religious and Community Resources}
\begin{itemize}
\item Scripture Earth: New Testament in print and audio, with Android apps: \url{https://www.scriptureearth.org/00eng.php?iso=poe}
\end{itemize}

\section{Digital and Audio Resources}
\begin{itemize}
\item Global Recordings Network, evangelistic and story audio for oral communities: \url{https://www.globalrecordings.net/en/language/poe}
\end{itemize}

\section{Curriculum Development Context}
The language is not known to be taught in schools. The available material is Scripture, short audio recordings and legacy primers in a superseded orthography. A modern orthography, a dictionary and a learner course were not found online, so curriculum work would depend on local speakers and on locating the 1995 reference works.

\end{document}